\chapter{パターンマッチ指向プログラミング問題集}

\section*{\texttt{member?}関数}

下記の空白をうめて\lstinline{member?}関数を完成させよ．

\begin{lstlisting}[language=egison]
member? x xs :=
  match xs as list eq with
  | `\fcolorbox{black}[gray]{1.0}{\vphantom{0}\hspace{5em}1\vphantom{0}\hspace{5em}}` -> True
  | _ -> False

member? 1 [1,2,3,4]
-- True

member? 5 [1,2,3,4]
-- False
\end{lstlisting}

\section*{\texttt{concat}関数}

下記の空白をうめて\lstinline{concat}関数を完成させよ．

\begin{lstlisting}[language=egison]
concat xss :=
  matchAllDFS xss as `\fcolorbox{black}[gray]{1.0}{\vphantom{0}\hspace{5em}1\vphantom{0}\hspace{5em}}` with
  | `\fcolorbox{black}[gray]{1.0}{\vphantom{0}\hspace{5em}2\vphantom{0}\hspace{5em}}` -> `\fcolorbox{black}[gray]{1.0}{\vphantom{0}\hspace{5em}3\vphantom{0}\hspace{5em}}`

concat [[1,2],[3],[4,5]]
-- [1,2,3,4,5]
\end{lstlisting}

\section*{四つ子素数}

四つ子素数を列挙するプログラムを記述せよ．
四つ子素数とは，$(p,p + 2,p + 6,p + 8)$という形の素数の四つ組のことをいう．

\begin{lstlisting}[language=egison]
take 4 (matchAll primes as list integer with
        | `\fcolorbox{black}[gray]{1.0}{\vphantom{0}\hspace{5em}1\vphantom{0}\hspace{5em}}` -> (p, p + 2, p + 6, p + 8))
-- [(5, 7, 11, 13), (11, 13, 17, 19), (101, 103, 107, 109), (191, 193, 197, 199)]
\end{lstlisting}

\section*{差が$6$の素数のペア}

下記の空白をうめて素数のペア$(p, p+6)$を抽出せよ．

\begin{lstlisting}[language=egison]
take 6 (matchAll primes as list integer with
        | `\fcolorbox{black}[gray]{1.0}{\vphantom{0}\hspace{5em}1\vphantom{0}\hspace{5em}}` -> (p, p + 6))
-- [(5, 11), (7, 13), (11, 17), (13, 19), (17, 23), (23, 29)]
\end{lstlisting}

\section*{\texttt{intersect}関数}

下記の空白をうめて，2つのコレクションの共通部分を抽出する\lstinline{intersect}関数を完成させよ．

\begin{lstlisting}[language=egison]
intersect xs ys :=
  matchAllDFS (xs, ys) as `\fcolorbox{black}[gray]{1.0}{\vphantom{0}\hspace{5em}1\vphantom{0}\hspace{5em}}` with
  | `\fcolorbox{black}[gray]{1.0}{\vphantom{0}\hspace{5em}2\vphantom{0}\hspace{5em}}` -> `\fcolorbox{black}[gray]{1.0}{\vphantom{0}\hspace{5em}3\vphantom{0}\hspace{5em}}`

intersect [1,2,3,4] [2,4,5]
-- [2,4]
\end{lstlisting}

\section*{\texttt{difference}関数}

下記の空白をうめて，第二引数のコレクションに含まれていない第一引数のコレクションの要素のみを抽出する\lstinline{difference}関数を完成させよ．

\begin{lstlisting}[language=egison]
difference xs ys :=
  matchAllDFS (xs, ys) as `\fcolorbox{black}[gray]{1.0}{\vphantom{0}\hspace{5em}1\vphantom{0}\hspace{5em}}` with
  | `\fcolorbox{black}[gray]{1.0}{\vphantom{0}\hspace{5em}2\vphantom{0}\hspace{5em}}` -> `\fcolorbox{black}[gray]{1.0}{\vphantom{0}\hspace{5em}3\vphantom{0}\hspace{5em}}`

difference [1,2,3,4] [2,4,5]
-- [1,3]
\end{lstlisting}

\section*{\texttt{doubles}関数}

コレクション中にちょうど$2$回現れる要素のみを抽出する\lstinline{doubles}関数を完成させよ．

\begin{lstlisting}[language=egison]
doubles xs :=
  matchAllDFS xs as `\fcolorbox{black}[gray]{1.0}{\vphantom{0}\hspace{5em}1\vphantom{0}\hspace{5em}}` with
  | `\fcolorbox{black}[gray]{1.0}{\vphantom{0}\hspace{5em}2\vphantom{0}\hspace{5em}}` -> `\fcolorbox{black}[gray]{1.0}{\vphantom{0}\hspace{5em}3\vphantom{0}\hspace{5em}}`

doubles [1,2,3,2,1,1,4,5,3]
-- [2,3]
\end{lstlisting}

\section*{\texttt{tails}関数}

\lstinline{tails}関数をパターンマッチを使って定義せよ．
ループ・パターンを使う回答とそうでない回答の2つを考えてみよ．

\begin{lstlisting}[language=egison]
tails xs :=
  matchAll xs as list something with
  | `\fcolorbox{black}[gray]{1.0}{\vphantom{0}\hspace{5em}1\vphantom{0}\hspace{5em}}`
  -> ts

tails [1,2,3,4]
-- [[1, 2, 3, 4], [2, 3, 4], [3, 4], [4], []]
\end{lstlisting}

\section*{多重集合の同値性チェック}

多重集合の同値性をチェックする2引数述語\lstinline{equalMultiset}をパターンマッチを使って定義してみよ．

\begin{lstlisting}[language=egison]
equalMultiset xs ys :=
  match (xs, ys) as (list eq, multiset eq) with
  | ([], [])
  -> True
  | `\fcolorbox{black}[gray]{1.0}{\vphantom{0}\hspace{5em}1\vphantom{0}\hspace{5em}}`
  -> equalMultiset xs' ys'
  | _
  -> False

equalMultiset [1,1,2] [2,1,1]
-- True

equalMultiset [1,1,2] [1,2,2]
-- False
\end{lstlisting}

ループ・パターンを使って再帰関数を使わずに書いてみよ．

\begin{lstlisting}[language=egison]
equalMultiset xs ys :=
  match (xs, ys) as (list eq, multiset eq) with
  | `\fcolorbox{black}[gray]{1.0}{\vphantom{0}\hspace{5em}2\vphantom{0}\hspace{5em}}`
  -> True
  | _
  -> False
\end{lstlisting}

\section*{ジョーカーの存在を考慮するポーカーの役判定}

\ref{pmo}章\ref{pmo-poker}節のポーカーの役判定のためのパターンマッチ（\lstinline{poker}関数）の実装を一切変更せずに，\lstinline{card}マッチャーの定義だけを変更して，ジョーカーの存在を考慮するポーカーの役判定をできるようにせよ．
\begin{lstlisting}[language=egison]
card := matcher
  | card $ $ as (suit, mod 13) with 
    | Card $s $n -> [(s, n)]
    | Joker -> `\fcolorbox{black}[gray]{1.0}{\vphantom{0}\hspace{5em}1\vphantom{0}\hspace{5em}}`
  | $ as something with
    | $tgt -> [tgt]

poker [Card Spade 5, Card Spade 6, Joker, Card Spade 8, Card Spade 9]
-- "Straight flush"
poker [Card Spade 5, Card Diamond 5, Joker, Card Club 5, Card Heart 7]
-- "Four of a kind"
\end{lstlisting}
